%!TEX root = ../../main.tex
%% ===============================================================
%% 7.8 도메인 지식 처치 T1--T7
%% 파일: chapters/07_models/08_treatments.tex
%% 상위: chapters/07_models/chapter.tex
%% 정의 라벨: sec:model-treatments
%% 참조 라벨: ch:localization
%% 포함 객체: -
%% 인용: appel2005technical, bahdanau2015neural, caruana1997multitask, kendall2018multi, lin2017focal, muller2019when, ruder2017overview, szegedy2016rethinking
%% 상태: 초안 (docs/OUTLINE.md 참고)
%% ===============================================================

\section{도메인 지식 처치 T1--T7}
\label{sec:model-treatments}

RQ5의 일곱 처치는 각각 LoL 도메인 지식과 학습 이론의 근거를 갖는 독립적인 설계 요소이며, 기본 구성에서 하나씩 켜서 효과를 측정한다.

\begin{description}[leftmargin=2.5em,style=nextline]
  \item[T1 초점 손실 \cite{lin2017focal}] $\mathcal{L}_{\mathrm{FL}} = -\alpha_t (1-p_t)^\gamma \log p_t$, $\gamma = 2$, $\alpha = 0.25$. 골드 격차가 큰 ``압도'' 교전은 쉬운 표본이므로 그 기여를 줄여 대등한 교전에 학습을 집중시킨다. $p_t = 0.9$인 표본의 가중치는 $0.01$이다.
  \item[T2 게임 국면 부호화] $\phi(t) = [\sigma((14-t)/\tau),\ \sigma((t-10)/\tau)\,\sigma((28-t)/\tau),\ \sigma((t-22)/\tau)]$, $\tau = 3$ 분. 초반(라인전)·중반(오브젝트 경합)·후반(5:5 한타)을 매끄러운 3 차원 벡터로 전역 특징에 추가한다.
  \item[T3 시간 주의 풀링 \cite{bahdanau2015neural}] $e_t = w^\top\tanh(W_a h_t + b_a)$, $\alpha_t = \softmax(e)_t$, $c = \sum_t \alpha_t h_t$, 출력 $[h_L \,\|\, c]$. 마지막 은닉 상태만 쓰는 대신 파워 스파이크 등 중요한 시점을 선택적으로 참조한다.
  \item[T4 모멘텀 특징] 금융 기술 분석의 MACD \cite{appel2005technical}에서 착안하여 1차 차분 $\Delta x_t$의 단기 평균($k = 3$)과 장기 평균의 차 $\mu_{\mathrm{short}} - \mu_{\mathrm{long}}$을 테이블 특징에 추가한다. 추세와 가속을 구분한다.
  \item[T5 역할 인식 인접행렬] $A'_{ij} = A_{ij}\,\softplus(R_{\mathrm{role}(i),\mathrm{role}(j)})$, $R \in \R^{5\times 5}$ 학습, 초기값 0($\softplus(0) = \ln 2$). 같은 거리라도 정글--미드 관계가 탑--서포터 관계보다 중요할 수 있음을 학습하게 한다.
  \item[T6 다중 과제 손실 \cite{caruana1997multitask,ruder2017overview}] $\mathcal{L} = \mathcal{L}_{\mathrm{fight}} + \lambda_g\|\hat y_{\mathrm{gold}} - y_{\mathrm{gold}}\|^2 + \lambda_k\|\hat y_{\mathrm{kill}} - y_{\mathrm{kill}}\|^2 + \lambda_o\|\hat y_{\mathrm{obj}} - y_{\mathrm{obj}}\|^2$, $(\lambda_g, \lambda_k, \lambda_o) = (0.1, 0.05, 0.05)$. 선택적으로 불확실성 가중 \cite{kendall2018multi}을 지원한다.
  \item[T7 라벨 평활화 \cite{szegedy2016rethinking}] $y_{\mathrm{smooth}} = y(1-\epsilon) + \epsilon/2$, $\epsilon = 0.05$. 균등 분포 방향의 KL 정규화로 과신을 줄인다 \cite{muller2019when}.
\end{description}

이 밖에 보간·검출기 절제로 T8(보간 없이 60 초 기준선), T9(킬 궤적 덮어쓰기 없는 검출기), T10(60 초 격자만 사용하는 검출기)을 등록하여 제 \ref{ch:localization} 장의 설계 요소를 같은 프로토콜로 평가할 수 있게 하였다.
